% ============================================================================
% APPENDIX B - DETAILED STATISTICAL TESTS
% ============================================================================

\labelsec{appendix_statistics}

\section{B.1 Correlation Significance}

\subsection{Vij-Mamton Pearson Correlation}

Sample: $N = 65$ structures with all three SHT units computed.

\begin{align}
    r &= 0.8285 \\
    t &= \frac{r \sqrt{N-2}}{\sqrt{1-r^2}} = \frac{0.8285 \times \sqrt{63}}{\sqrt{1-0.6864}} \\
    &= \frac{0.8285 \times 7.9373}{0.5601} = \frac{6.5782}{0.5601} = \mathbf{11.744} \\
    p &= 1.633 \times 10^{-17} \quad \text{(two-tailed, df=63)}
\end{align}

\textbf{Interpretation:} The probability of observing $|r| \geq 0.83$ by chance in a dataset of 65 pairs is $1.6 \times 10^{-17}$.

\subsection{95\% Confidence Interval (Fisher z-transformation)}

\begin{align}
    z &= \frac{1}{2}\ln\!\left(\frac{1+r}{1-r}\right) = \frac{1}{2}\ln\!\left(\frac{1.8285}{0.1715}\right) = \frac{1}{2}\ln(10.662) = 1.1764 \\
    \text{SE}(z) &= \frac{1}{\sqrt{N-3}} = \frac{1}{\sqrt{62}} = 0.12700 \\
    z_{95\%} &= 1.1764 \pm 1.9600 \times 0.12700 = [0.9276,\;1.4252] \\
    r_{\rm lo} &= \tanh(0.9276) = 0.7327 \\
    r_{\rm hi} &= \tanh(1.4252) = 0.8920 \\
    \mathbf{95\%\,\text{CI}} &= [0.733,\;0.892]
\end{align}

All values in the interval exceed $r = 0.4$ (the threshold for a ``large'' correlation per Cohen's conventions). The lower bound alone ($r = 0.73$) would be the strongest inter-unit correlation reported in DNA structural biology.

\subsection{Vij-Shambhian and Shambhian-Mamton}

\begin{align}
    r(v_0, \eta_r) &= -0.0851, \quad t = -0.679, \quad p = 0.500 \quad \text{(NS)} \\
    r(\eta_r, \Omega) &= -0.1522, \quad t = -1.224, \quad p = 0.226 \quad \text{(NS)}
\end{align}

Neither is significant at any conventional threshold. Shambhian is statistically independent of both Vij and Mamton.

\section{B.2 Mann-Whitney U Tests for A-tract Validation}

All tests two-tailed. Groups: A-tract ($n_A = 10$), Control ($n_B = 15$).

\subsection{Vij Separation}

\begin{align}
    \bar{v}_{0,A} &= 0.0797 \pm 0.012 \quad \bar{v}_{0,B} = 0.1238 \pm 0.031 \\
    U &= 21, \quad p = 1.79 \times 10^{-8}
\end{align}

\subsection{Shambhian Separation}

\begin{align}
    \bar{\eta}_{r,A} &= 0.9928 \pm 0.005 \quad \bar{\eta}_{r,B} = 0.9650 \pm 0.015 \\
    p &= 6.08 \times 10^{-34}
\end{align}

This is the smallest p-value in the study. The directness of the Shambhian A-tract effect reflects the tight minor groove geometry in A-tract DNA: propeller-twisted A$\cdot$T pairs produce a characteristic cross-sectional eccentricity pattern.

\subsection{Mamton Separation}

\begin{align}
    \bar{\Omega}_A &= 70.5 \quad \bar{\Omega}_B = 3{,}438.9 \\
    p &\approx 0 \quad \text{(below double-precision float)}
\end{align}

The exact p-value cannot be represented as a 64-bit float; it is more extreme than $\sim 5 \times 10^{-324}$.

\section{B.3 Effect Sizes (Cohen's d)}

\begin{table}[h]
    \centering
    \caption{Effect Sizes for A-tract vs Control Separation}\labeltab{effect_sizes}
    \begin{tabular}{lll}
        \toprule
        \textbf{Unit} & \textbf{Cohen's $d$} & \textbf{Interpretation} \\
        \midrule
        Vij       & $d = 1.93$ & Very large ($d > 0.8$: large) \\
        Shambhian & $d = 2.70$ & Very large \\
        Mamton    & $d > 3.0$  & Extreme (skewed distribution; non-parametric preferred) \\
        \bottomrule
    \end{tabular}
\end{table}

\section{B.4 Z-Score Outlier Table}

Computed from the merged N=65 dataset. Z-scores = (value $-$ mean) / std.

\begin{table}[h]
    \centering
    \caption{Z-scores for Selected Notable Structures}\labeltab{z_scores}
    \begin{tabular}{lrrrrrr}
        \toprule
        \textbf{PDB} & $v_0$ & $\eta_r$ & $\Omega$ & $z_{v_0}$ & $z_{\eta_r}$ & $z_\Omega$ \\
        \midrule
        1AOI & 0.311 & 0.988 & 26,291 & $+4.73$ & $+0.54$ & $+5.55$ \\
        1F66 & 0.297 & 0.828 & 26,441 & $+4.41$ & $-2.18$ & $+5.58$ \\
        1AHD & 0.088 & 0.710 & 94.4   & $-0.39$ & $-4.18$ & $-0.19$ \\
        1CQT & 0.174 & 0.986 & 1,049  & $+1.59$ & $+0.51$ & $+0.02$ \\
        1EYG & 0.115 & 0.869 & 1,427  & $+0.23$ & $-1.49$ & $+0.10$ \\
        1LO1 & ---   & 0.536 & ---    & ---     & $-6.2$  & ---     \\
        1I8M & ---   & 0.9998 & ---   & ---     & $\sim+0.6$ & ---  \\
        \bottomrule
    \end{tabular}
\end{table}

\section{B.5 Null Hypothesis Test Results}

Permutation test: 10,000 random permutations of structure labels. Real dataset deviates by:
\begin{align}
    Z_{\rm CORR} &= -40.22 \quad \text{(correlation structure far below null)} \\
    Z_{\rm VAR} &= +56.98 \quad \text{(variance structure far above null)}
\end{align}

Both $|Z| > 5$, corresponding to $p < 3 \times 10^{-7}$ per test. The SHT signal is not a numerical artifact.

\section{B.6 Classification Power Analysis}

For the discriminant between nucleosome and free-duplex classes (2 vs 16 structures):

Using the Mamton effect size $d > 3.0$, expected power at $N_{\rm nucl} = 2$ is already $> 0.99$ (two-sample t-test approximation). The small nucleosome sample is sufficient because the effect size is extreme.

For the full 4-class classifier at $N = 33$, the 5-fold CV with $\pm 11.7\%$ uncertainty is consistent with bootstrapped confidence intervals; 91\% accuracy falls in the 95\% CI $[\approx 72\%, \approx 100\%]$.

\section{B.7 Linear Model Coefficients and Residuals}

Linear fit $\Omega = \beta_1 v_0 + \beta_0$ (OLS, $N = 65$):

\begin{align}
    \hat{\beta}_1 &= 86{,}683.4 \quad \text{(slope)} \\
    \hat{\beta}_0 &= -8{,}109.9 \quad \text{(intercept)} \\
    R^2 &= 0.686 \quad \text{(69\% variance explained)} \\
    F(1,63) &= 137.9, \quad p = 1.6 \times 10^{-17}
\end{align}

Residuals are heteroscedastic (larger at high $\Omega$) consistent with the log-normal distribution of Mamton. A log-linear model ($\ln\Omega = \beta_1 v_0 + \beta_0$) gives $R^2 = 0.59$ with more uniform residuals.

